\documentclass[UTF8,12pt]{ctexart}
\usepackage[paperwidth=240mm,paperheight=200mm,margin=14mm,top=8mm]{geometry}
\usepackage{amsmath,amssymb,mathtools}
\usepackage{xcolor}
\pagestyle{empty}
\setlength{\parindent}{0pt}
\setlength{\parskip}{0pt}
\newcommand{\qn}[1]{\textbf{\large #1}\ }
\xeCJKDeclareCharClass{CJK}{"2460 -> "24FF}
\begin{document}
\qn{1994.1}
若 $f(x)=\begin{cases}\dfrac{\sin 2x+\mathrm{e}^{2ax}-1}{x}, & x\neq 0,\\ a, & x=0\end{cases}$ 在 $(-\infty,+\infty)$ 上连续，则 $a=$______.

\vfill
\newpage
\qn{1994.2}
设函数 $y=y(x)$ 由参数方程 $\begin{cases}x=t-\ln(1+t),\\ y=t^3+t^2\end{cases}$ 所确定，则 $\dfrac{\mathrm{d}^2y}{\mathrm{d}x^2}=$______.

\vfill
\newpage
\qn{1994.3}
$\dfrac{\mathrm{d}}{\mathrm{d}x}\left(\int_0^{\cos 3x}f(t)\,\mathrm{d}t\right)=$______.

\vfill
\newpage
\qn{1994.4}
$\int x^3\mathrm{e}^{x^2}\,\mathrm{d}x=$______.

\vfill
\newpage
\qn{1994.5}
微分方程 $y\,\mathrm{d}x+(x^2-4x)\,\mathrm{d}y=0$ 的通解为______.

\vfill
\newpage
\qn{1994.6}
设 $\lim\limits_{x\to0}\dfrac{\ln(1+x)-(ax+bx^2)}{x^2}=2$，则（　　）
\\
\noindent (A) $a=1$，$b=-\dfrac{5}{2}$.　　(B) $a=0$，$b=-2$.
\\
\noindent (C) $a=0$，$b=-\dfrac{5}{2}$.　　(D) $a=1$，$b=-2$.

\vfill
\newpage
\qn{1994.7}
设 $f(x)=\begin{cases}\dfrac{2}{3}x^3, & x\leqslant 1,\\ x^2, & x>1,\end{cases}$ 则 $f(x)$ 在点 $x=1$ 处的（　　）
\\
\noindent (A) 左、右导数都存在.　　(B) 左导数存在，但右导数不存在.
\\
\noindent (C) 左导数不存在，但右导数存在.　　(D) 左、右导数都不存在.

\vfill
\newpage
\qn{1994.8}
设 $y=f(x)$ 是满足微分方程 $y''+y'-\mathrm{e}^{\sin x}=0$ 的解，且 $f'(x_0)=0$，则 $f(x)$ 在（　　）
\\
\noindent (A) $x_0$ 的某个邻域内单调增加.　　(B) $x_0$ 的某个邻域内单调减少.
\\
\noindent (C) $x_0$ 处取得极小值.　　(D) $x_0$ 处取得极大值.

\vfill
\newpage
\qn{1994.9}
曲线 $y=\mathrm{e}^{\frac{1}{x^2}}\arctan\dfrac{x^2+x+1}{(x-1)(x+2)}$ 的渐近线有（　　）
\\
\noindent (A) 1条.　　(B) 2条.
\\
\noindent (C) 3条.　　(D) 4条.

\vfill
\newpage
\qn{1994.10}
设 $M=\int_{-\frac{\pi}{2}}^{\frac{\pi}{2}}\dfrac{\sin x}{1+x^2}\cos^4 x\,\mathrm{d}x$，$N=\int_{-\frac{\pi}{2}}^{\frac{\pi}{2}}(\sin^3 x+\cos^4 x)\,\mathrm{d}x$，$P=\int_{-\frac{\pi}{2}}^{\frac{\pi}{2}}(x^2\sin^3 x-\cos^4 x)\,\mathrm{d}x$，则有（　　）
\\
\noindent (A) $N<P<M$.　　(B) $M<P<N$.
\\
\noindent (C) $N<M<P$.　　(D) $P<M<N$.

\vfill
\newpage
\qn{1994.11}
设 $y=f(x+y)$，其中 $f$ 具有二阶导数，且其一阶导数不等于 $1$，求 $\dfrac{\mathrm{d}^2y}{\mathrm{d}x^2}$.

\vfill
\newpage
\qn{1994.12}
计算 $\int_0^1 x(1-x^4)^{\frac{3}{2}}\,\mathrm{d}x$.

\vfill
\newpage
\qn{1994.13}
计算 $\lim\limits_{n\to\infty}\tan^n\left(\dfrac{\pi}{4}+\dfrac{2}{n}\right)$.

\vfill
\newpage
\qn{1994.14}
计算 $\int\dfrac{\mathrm{d}x}{\sin 2x+2\sin x}$.

\vfill
\newpage
\qn{1994.15}
如图，设曲线方程为 $y=x^2+\dfrac{1}{2}$，梯形 $OABC$ 的面积为 $D$，曲边梯形 $OABC$ 的面积为 $D_1$，点 $A$ 的坐标为 $(a,0)$，$a>0$. 证明：$\dfrac{D}{D_1}<\dfrac{3}{2}$.
设当 $x>0$ 时，方程 $kx+\dfrac{1}{x^2}=1$ 有且仅有一个解，求 $k$ 的取值范围.
设 $y=\dfrac{x^3+4}{x^2}$，

\vfill
\newpage
\qn{1994.16}
求函数的增减区间及极值；

\vfill
\newpage
\qn{1994.17}
求函数图形的凹凸区间及拐点；

\vfill
\newpage
\qn{1994.18}
求其渐近线；

\vfill
\newpage
\qn{1994.19}
作出其图形.
求微分方程 $y''+a^2y=\sin x$ 的通解，其中常数 $a>0$.
设 $f(x)$ 在 $[0,1]$ 上连续且递减，证明：当 $0<\lambda<1$ 时，$\int_0^\lambda f(x)\,\mathrm{d}x\geqslant\lambda\int_0^1 f(x)\,\mathrm{d}x$.
求曲线 $y=3-|x^2-1|$ 与 $x$ 轴围成的封闭图形绕直线 $y=3$ 旋转所得的旋转体体积.

\vfill
\newpage
\end{document}